\section{Overview of our proof}  \label{sec:overview}

In this section we give a more detailed overview of the technical
parts of the paper.

\subsection{Basic quantum notation and qudits}
While the main body of the paper contains a more complete set of
quantum preliminaries in~\Cref{sec:prelims}, for the purposes of this
introduction we define some basic notation, aimed at the reader who is
familiar with the standard quantum computing formalism over qubits but
is less familiar with \emph{qudits}: quantum systems of dimension not equal to $2$.
In this paper, we make extensive use of such qudits: in particular, for a finite field $\F_Q$, we will
consider qudits of dimension $Q$, with a basis state $\ket{i}$ for
every element $i \in \F_Q$. Under tensor product, we obtain a basis
for the space of $M$ qudits of dimension $Q$ where each basis state
$\ket{x}$ corresponds to a vector $x \in \F_Q^M$.

The basic resource state used in our protocols will be the EPR state
over $2M$ qudits of dimension $Q$. The qudits are split into two
registers of $M$ qudits each, held by the two provers Alice and Bob,
respectively.
\[ \ket{\epr_Q^M} = \frac{1}{\sqrt{Q^M}} \sum_{x \in \F_Q^M} \ket{x}_{\reg{Alice}}
  \ot \ket{x}_{\reg{Bob}}. \]
This state is a \emph{maximally entangled} state between Alice and
Bob.

Acting on this state, we will ask the provers to perform measurements
from a special class called \emph{Pauli basis measurements}. To define
these over a general field $\F_Q$ requires the introduction of some
finite field technology, in particular the finite field trace
function. For simplicity, in this overview we will imagine that $Q$ is
prime, allowing the addition in $\F_Q$ to be identified with the
additive group $\mathbb{Z}_Q$, and simplifying the definition of the
Paulis; in the main body of the paper, we will work with $Q$ a power
of $2$. For a single qudit of dimension $Q$, the Pauli $X$ and $Z$
bases are the sets $\{\ket{\tau^X_u}\}_{u \in \F_Q}$ and $\{\ket{\tau^Z_u}\}_{u
  \in \F_Q}$ of vectors
\[ \ket{\tau^X_u} = \frac{1}{\sqrt{Q}}\sum_{x \in \F_Q}  \omega^{x u} \ket{x}, \qquad
  \ket{\tau^Z_u} = \ket{u}, \]
where $\omega = \exp(2\pi i/Q)$ is the $Q$-th root of unity. We denote
the projectors onto these basis states by $\tau^X_u$ and $\tau^Z_u$,
respectively. For a system of $M$ qudits, the Pauli
$X$ and $Z$ observables are a set of \emph{generalized observables}
indexed by elements of $\F_Q^M$: a generalized observable is a
Hermitian matrix with eigenvalues that are $Q$-th roots of unity.
They are given by
\[ X(v) =\sum_{u \in \F_Q^M} \omega^{u \cdot v} \tau^X_{u_1} \ot \dots
  \ot \tau^Z_{u_M}, 
  \qquad Z(v) = \sum_{u \in \F_Q^M} \omega^{u \cdot v} \tau^Z_{u_1}
  \ot \dots \ot \tau^Z_{u_M},\]
where $u_1, \dots, u_M$ are the components of the vector $u$, and
$u\cdot v$ is the dot product $\sum_{i=1}^{M} u_i \cdot
v_i$. Measuring a generalized observable means performing a projective
measurement onto the eigenvectors of the observable, with the
outcome $a$ corresponding to the eigenvector with eigenvalue
$\omega^{a}$. 



\subsection{Our starting point: a classical interactive proof for
  $\NEEXP$}

We start with a classical multiprover interactive proof protocol for
$\NEEXP$. The equality $\MIP = \NEXP$ was originally shown by Babai,
Fortnow, and Lund~\cite{BFL91} using a protocol based on the \emph{multilinearity
  test}: the idea is that an exponentially-long witness for a problem
in $\NEXP$ is encoded in the truth-table of a multivariate polynomial
function over a finite field, which is linear in each of the variables
individually. The verifier is able to verify the witness
by evaluating the multilinear polynomial over appropriately chosen
points and subspaces. To scale up to $\NEEXP$, we use a much more
efficient version of the same idea, replacing the multilinearity test
with the \emph{low-degree test}, which works with multivariate
polynomials of low total degree. This more efficient construction
comes from the PCP literature. We give a relatively self-contained
presentation of the protocol in~\Cref{sec:classical-pcp}. For the
purposes of this overview, it is sufficient to know the following: any
problem in $\NEEXP$ can be reduced to satisfiability for a doubly exponentially long 3Sat
formula, succinctly encoded by a polynomial-sized circuit. (We refer
to this problem as $\succinctsquared$). Given a 3Sat formula
$\psi$, we would like the provers to prove to us that they have a
satisfying assignment $a$ to this formula. Instead of reading the
assignment directly, we will ask the provers to
encode their assignment as a multivariate polynomial $g_a: \F_Q^M \to
\F_Q$, where the number of variables $M$ and the finite field size $Q$
are appropriately chosen parameters, and return evaluations of this polynomial. To check that $a$
satisfies $\psi$, the verifier first uses a technique called
arithmetization to convert the
formula $\psi$ into a multivariate polynomial $g_{\psi}: \F_Q^{3M + k}
\to \F_Q$. The polynomial $g_{\psi}$ is chosen such that the assignment $a$
satisfies $\psi$ if and only if the expression
\[ \mathrm{sat}_{\psi,a}(x,b, w) := g_{\psi}(x,b,w) \cdot (g_a(x_1)
  -b_1)(g_a(x_2) - b_2)(g_a(x_3) - b_3) \]
is equal to $0$ at every point in a particular subset $H \subseteq \F_Q^{3M
  +k}$. Our classical protocol for $\NEEXP$ checks this condition:
\begin{infthm}[\Cref{sec:classical-pcp} in the body]
There exists a protocol $\game_0$ for $\mathsf{Succinct}$-$\mathsf{Succinct}$-$\mathsf{3Sat}$ (and hence
$\NEEXP$), where the verifier's questions to the provers are
constant-dimension subspaces of $\F_Q^{M}$, and the provers' responses
are evaluations of degree-$D$ $M$-variate polynomials on these
subspaces. The parameters $M, Q, D$ are all chosen to be $\exp(n)$,
and hence the question and answer lengths as well as the runtime of
the verifier in this protocol are
$\exp(n)$. 
\end{infthm}
The distribution over subspaces sent to the provers in $\game_0$ is
relatively simple, and in fact is independent of the instance of
$\succinctsquared$ being tested. For the purposes of this overview, the
reader can take the distribution over pairs of questions to be the
\emph{plane-point distribution} $\calD$. A pair $(\bs, \bu) \sim
\calD$ consists of a uniformly random affine plane $\bs \subseteq
\F_Q^M$, which is sent to Alice, and a uniformly random point $\bu
\in \bs$ which is sent to Bob. The full distribution over
questions in $\game_0$ is more complicated than this but the essential
ideas of our protocol will be illustrated by restricting to this case.
\subsection{Restricting the strategies: registers and compilers}
One of the main challenges in working with entangled provers is
showing soundness against general entangled strategies. An important
technique in this area is to force the provers to use a particular
state and class of measurements by playing a type of game known as a
\emph{self-test}.
\begin{infdef}
  A game $\game_{\mathrm{test}}$ is a \emph{self-test} for a state
  $\ket{\psi}$ and measurements $M^x$ if any strategy that succeeds in
  $\game_{\mathrm{test}}$ with probability $1- \eps$ must use a state
  $\ket{\psi'}$ and measurements $(M')^x$ that are
  $\delta(\eps)$-close, in the appropriate metric, to $\ket{\psi}$ and $M^x$.
\end{infdef}
Some of the earliest self-tests include the famous CHSH game, which
self-tests the Pauli $X$ and $Z$ operators on a single EPR pair (of
qubits). Self-testing technology has greatly advanced over the years,
and in this paper we design a highly efficient self-test based on the
low-degree test of~\cite{NV18a}.
\begin{infthm}[\Cref{thm:basis-test} in the body]
  The Pauli basis test $\paulistrat{n}{q}$ is a self-test for the
  state $\ket{\epr_q^n}$ and the Pauli $X$ and $Z$ basis
  measurements. This test sends the players questions of length
  $O(\log(n))$ and receives answers of length $O(\poly(n))$.
\end{infthm}
The Pauli $X$ and $Z$ measurements are ``complete'' measurements, and
as a consequence, there is no nontrivial measurement on a set $n$
qudits that can be measured jointly with both the Pauli $X$ and $Z$
measurements on those qudits. Using this property, we design a game
called the \emph{data-hiding game}, which certifies that a prover's
measurements act trivially on a specified set of qudits.
\begin{infthm}[\Cref{thm:data-hiding-game} in the body]
  The data-hiding game $\game_{\mathrm{hide}}$ is a self test for
  states $\ket{\psi} = \ket{\epr_q^n} \ot \ket{\rmaux}$ and
  measurements $M^x$ of the form $M^x = I \ot (M')^x_{\reg{aux}}$. It
  has questions of length $O(\log(n))$ and answers of length $O(\poly(n))$.
\end{infthm}
Together, the Pauli basis test and the data-hiding game allow us to restrict our analysis of our protocols to a
class of strategies we call \emph{register strategies}: strategies for
which the shared state is a collection of $\ell$ registers, each in an
EPR state, together with some auxiliary register:
\[ \ket{\psi} = \ket{\epr_{q_1}^{n_1}} \ot \dots \ot
  \ket{\epr_{q_\ell}^{n_\ell}} \ot \ket{\rmaux}, \]
and where the provers can be commanded to perform either (1) Pauli basis measurements on specified subsets of the
registers, or (2) measurements that do \emph{not} act on specified
subset of the EPR registers (but act on the auxiliary register or the
remaining EPR registers). We formalize this by designing a \emph{compiler},
which takes in a protocol $\game$ that is complete and sound for register
strategies, and produces a new protocol $\game'$ which is complete and
sound over all strategies.
\begin{infthm}[\Cref{theorem:one-to-zero-compiler}
  and~\Cref{theorem:two-to-one-compiler} in the body]\label{thm:compileroonie}
  Suppose $\game$ is a protocol for a computation problem for which completeness and soundness
  hold for register strategies, with $O(1)$ many registers of size
  $n$. (That is, for YES instances of the problem, there exists a
  register strategy achieving value $1$, and for NO instances, no
  register strategy achieves value greater than $1/2$). Let the
  questions in $\game$ be of length $Q$ and the answers be of length
  $A$. Then there exists a protocol $\game'$ which is complete and
  sound for general strategies, and for which the question length is $Q
  + \log(n)$ and the answer length is $A + \poly(n)$.
\end{infthm}
The compiled protocol $\game'$ either runs the original protocol
$\game$, or, with some probability, runs the Pauli basis test, the
data-hiding game, or a consistency test.


\subsection{Question reduction through introspection}

With our compiler in place, we have now given the verifier the power
to command the provers to perform Pauli basis measurements on a set of
EPR pairs.
We would like to use this to reduce the question size of the classical protocol $\game_0$ for $\NEEXP$ described above
from $\exp(n)$ to $\poly(n)$.
We will do so by forcing the provers, rather than the verifier, to sample the protocol's $\exp(n)$-length questions,
a technique we call ``introspection".
That is, we would like to force the provers to
sample pairs $(\bs, \bu)$ from the plane-vs-point distribution $\calD$, where $\bs$ is a uniformly random
affine plane in $\F_Q^M$, and $\bu$ a uniformly random point on
$\bs$.

To design a scheme to sample from this distribution, let us first fix
a representation of affine planes. We will represent an affine plane
by an \emph{intercept} $u \in \F_Q^M$ and two \emph{slopes} $v_1, v_2
\in \F_Q^M$. The plane given by $u, v_1, v_2$ is the set $s_u^v = \{ u +
\lambda_1 v_1 + \lambda_2 v_2 : \lambda_1, \lambda_2 \in \F_Q\}$. As a
first attempt, we may try the following scheme:
\begin{enumerate}
  \item Alice and Bob share three registers, each of which contains an
    EPR state, so their shared state is
    \[ \ket{\psi_0} = \ket{\epr_Q^M}_{R_0} \ot \ket{\epr_Q^M}_{R_1} \ot
      \ket{\epr_Q^M}_{R_2}. \]
  \item Alice first measures her half of registers $R_1$ and $R_2$ in
    the Pauli $Z$-basis, to obtain uniformly random outcomes $\bv_1,
    \bv_2$. The shared state is now
    \[ \ket{\psi_1} = \ket{\epr_Q^M}_{R_0} \ot
      (\ket{\bv_1}_{\reg{Alice}} \ot \ket{\bv_1}_{\reg{Bob}})_{R_1} \ot
      (\ket{\bv_2}_{\reg{Alice}} \ot \ket{\bv_2}_{\reg{Bob}})_{R_2}. \]
  \item Now, Alice and Bob both measure register $R_0$ in the Pauli
    $Z$-basis, both obtaining the same outcome $\bu$. The shared state
    is now
    \[ \ket{\psi_2} = (\ket{\bu}_{\reg{Alice}} \ot
      \ket{\bu}_{\reg{Bob}})_{R_0}    \ot  (\ket{\bv_1}_{\reg{Alice}} \ot \ket{\bv_1}_{\reg{Bob}})_{R_1} \ot
      (\ket{\bv_2}_{\reg{Alice}} \ot \ket{\bv_2}_{\reg{Bob}})_{R_2}. \]
    Alice sets her plane $\bs$ to be $s_{\bu}^{\bv}$ and Bob sets his point to be $\bu$.
  \end{enumerate}
  Indeed, the pair $(\bs, \bu)$ generated by this procedure is
  distributed according to $\calD$. However, there is a problem:
  through her measurement, Alice obtains additional side information,
  specifically the value of Bob's point $\bu$. Can we command Alice to
  erase the side information? In fact, we can, using the
  \emph{Heisenberg uncertainty principle}: if two observables
  anticommute, then measuring one completely destroys information
  about the other. Using this idea, we modify our protocol as follows:
  \begin{enumerate}
  \item As above.
  \item As above. At this point,  applying the definition of
    $\ket{\epr_Q^M}$, we can write the shared state as
    \[ \ket{\psi_1} \propto \sum_{u \in \F_Q^M} (\ket{u}_{\reg{Alice}}
      \ot \ket{u}_{\reg{Bob}})_{R_0} \ot
      (\ket{\bv_1}_{\reg{Alice}} \ot \ket{\bv_1}_{\reg{Bob}})_{R_1} \ot
      (\ket{\bv_2}_{\reg{Alice}} \ot \ket{\bv_2}_{\reg{Bob}})_{R_2}. \]
  \item {\bf New:} Intuitively, we would like Alice to be
    \emph{prevented} from measuring the component of the intercept
    along the directions $\bv_1, \bv_2$. This information would be
    obtained by measuring the observables\footnote{Strictly speaking,
      this is only true when $\bv_1 \cdot \bv_1 \neq 0$ and $\bv_2
      \cdot \bv_2 \neq 0$. A more rigorous treatment of this is given
      in~\Cref{sec:rotated-data-hiding}.} $Z(\bv_1), Z(\bv_2)$. To
    destroy it, we will ask Alice to measure the \emph{complementary}
    Pauli observables $X(\bv_1),
    X(\bv_2)$ on register $R_0$, obtaining outcomes $\balpha_1, \balpha_2
    \in \F_Q$. The shared state is now
    \begin{align*}
      \ket{\psi_2'} &\propto \sum_{u} \sum_{\lambda, \mu} \left(
                      \omega^{\balpha_1 \lambda + \balpha_2 \mu}
          \ket{\underbrace{u + \lambda \bv_1 + \mu \bv_2}_{u'}}_{\reg{Alice}}
        \ket{u}_{\reg{Bob}}\right)_{R_0} 
            (\ket{\bv_1}_{\reg{Alice}} \ot
                      \ket{\bv_1}_{\reg{Bob}})_{R_1} \\
      &\qquad\qquad\qquad\qquad\ot
        (\ket{\bv_2}_{\reg{Alice}} \ot \ket{\bv_2}_{\reg{Bob}})_{R_2}.
    \end{align*}
    where, as above, $\omega = \exp(2\pi i / Q)$ is a $Q$-th root of
    unity. Alice and Bob's state on $R_0$ is now a uniform 
    superposition over pairs $u, u'$ of points lying on the same
    affine subspace with slopes $\bv_1, \bv_2$.
  \item Alice and Bob both measure register $R_0$ in the $Z$ basis,
    obtaining outcomes $\bu$ and $\bu'$, respectively. The shared
    state is now
        \[ \ket{\psi_3'} = (\ket{\bu}_{\reg{Alice}} \ot
      \ket{\bu'}_{\reg{Bob}})_{R_0}    \ot  (\ket{\bv_1}_{\reg{Alice}} \ot \ket{\bv_1}_{\reg{Bob}})_{R_1} \ot
      (\ket{\bv_2}_{\reg{Alice}} \ot \ket{\bv_2}_{\reg{Bob}})_{R_2}. \]
Alice sets her
    plane to be $s^{\bv}_{\bu}$ and Bob sets his point to be
    $\bu'$. 
  \end{enumerate}
  Now, from the calculation performed above, it's clear that Bob's
  point $\bu'$ is uncorrelated with Alice's intercept $\bu$, apart
  from lying in the plane $s^{\bv}_{\bu}$, and hence there is no
  further information about Bob's point that Alice can learn by
  measuring her portion of the final state $\ket{\psi_3'}$. But Alice still obtains
  some additional information from her measurements along the way, in particular the
  outcomes $\alpha_1, \alpha_2$ of the $X$ measurements. And moreover,
  how can we certify that the $X$ measurements were performed
  correctly, since they are not Pauli basis measurements as given to
  us by the compiler? To answer these questions, we define a new game
  called the \emph{partial data-hiding game} (\Cref{thm:partial-data-hiding-game}), which certifies that
  Alice and Bob perform the steps described above and that no
  extra information is leaked. Building on this game, we can now
  design a protocol for $\NEEXP$ with small question size:

 \begin{infthm}[\Cref{thm:main-result-of-this-part} in the body]\label{thm:small-quest}
There is an $\MIP^*$ protocol $\game_{1}$ for $\NEEXP$ with questions
of length $\poly(n)$, and answers of length $\exp(n)$. The verifier can
generate the questions in $\poly(n)$ time but needs $\exp(n)$ time to
verify the answers.
\end{infthm}

\subsection{Answer reduction through PCP composition}
We have succeeded in obtaining a game with short questions, but the
answers are now exponentially long. In the last step, we will use
composition with a classical probabilistically checkable proof (PCP)
to delegate verification of the answers to the provers.

Schematically, the protocol $\game_{1}$ consists of the following
steps:
\begin{enumerate}
  \item The verifier sends Alice a question $\bx$ and Bob a question
    $\by$.
  \item Alice returns an (exponentially-long) answer $\bA$ and Bob an
    exponetially-long answer $\bB$.
  \item The verifier computes a verification predicate $V(\bx, \by,
    \bA, \bB)$ in exponential time.
  \end{enumerate}
  We would like to delegate the last step to the provers by asking
  them to compute a PCP proof that $V(\bx, \by,\bA,\bB) = 1$, which
  the verifier can check by communicating only polynomially many bits
  with the provers. However, we face
  an obstacle: Alice cannot know $\by$ and $\bB$, and neither can Bob
  know $\bx$ and $\bA$, and distributed PCPs (where neither party
  knows the entire assignment) are known to be
  impossible~\cite{ARW17}. To proceed, we will first have to modify
  $\game_{1}$ by \emph{oracularizing} it:
  \begin{enumerate}
    \item The verifier sends Alice the questions $\bx, \by$, and Bob
      either $\bx$ or $\by$, chosen uniformly at random.
    \item Alice returns exponentially-long answers $\bA, \bB$, and
      Bob returns an answer $\bC$.
    \item The verifier computes a verification predicate $V(\bx,
      \by, \bA, \bB)$ on Alice's questions and answers, and further
      checks that $\bA = \bC$, if Bob received $\bx$, or that $\bB =
      \bC$, if Bob received $\by$.
    \end{enumerate}
    The idea is that the new Alice simulates both Alice and Bob from
    the original protocol, and the new Bob certifies that the new
    Alice does not take advantage of her access to both questions to cheat.
    It is well-known that oracularization does not harm the soundness
    of interactive protocols, be they classical or quantum. However,
    in the quantum world, it is not necessarily the case that the
    oracularized protocol retains \emph{completeness}. This is because
    Alice and Bob may have been asked to perform non-compatible
    measurements in the original protocol, rendering it impossible for
    the new Alice to simulate both the original Alice and
    Bob. Fortunately for us, the honest strategy for protocol
    $\game_1$ is such that completeness under oracularization.
    
    Now that a single prover is in possession of all inputs to the
    verification predicate $V$, we can implement our idea of using a
    PCP proof. Classically, this idea is known as PCP
    \emph{composition}, and is extensively used in the PCP
    literature. In the quantum case, the requirement to maintain soundness against
    entanglement makes composition technically difficult, and we defer
    the details to \Cref{part:answer} of the paper. Once the
    composition is performed, we reach our main result.
    \begin{infthm}[\Cref{thm:main} in the body]\label{thm:whatevewre}
      There is an $\MIP^*$ protocol $\game_{2}$ for $\mathsf{Succinct}$-$\mathsf{Succinct}$-$\mathsf{3Sat}$
      (and hence for $\NEEXP$) with
      question size, answer, and verifier runtime $\poly(n)$.
    \end{infthm}
    
\subsection{Organization}

The paper is organized into five parts.  The first part is the introduction and this overview. The remaining parts are organized as follows.
\begin{itemize}
\item \Cref{part:prelims} contains two sections of preliminaries, one containing the classical background and another the quantum background.
\item \Cref{part:stack} contains the register compiler, i.e.\ the proof of \Cref{thm:compileroonie}.
	This involves designing the Pauli basis test (\Cref{sec:self-test-pauli}) and the data hiding test (\Cref{sec:data-hiding-layer}).
	\Cref{sec:register-overview} serves as an introduction to this part and contains more details on the organization.
\item \Cref{part:neexp} contains the ``introspection" question reduction step, i.e.\ the proof of \Cref{thm:small-quest}.
	To begin, we sketch the classical $\mip$ protocol for $\succinct$ in \Cref{sec:classical-pcp}.
	Then we give the introspected, i.e.\ ``big", low-degree test in \Cref{sec:big-degree},
	and finish by giving the entire small-question $\neexp$ protocol in \Cref{sec:big-neexp-protocol}.
	\Cref{sec:intersection} contains a test necessary for the protocol called the ``intersecting lines test".
	It allows us carry over the results of the low-degree test from one register to another.
\item \Cref{part:answer} contains the answer reduction, i.e.\ the proof of \Cref{thm:whatevewre}.
	The construction involves composing PCP protocols with error-correcting codes, and so \Cref{sec:error} surveys the properties we need of an error-correcting code.
	Finally, \Cref{sec:answer-reduction} contains the actual proof of the answer reduction step.
\end{itemize}

%%% Local Variables:
%%% mode: latex
%%% TeX-master: "main"
%%% End:
%%%---------- close: overview






\part{Preliminaries}

\label{part:prelims}

%%%%%%%%%%% jump to prelims-classical
%%%---------- open: prelims-classical
%!TEX root = main.tex
\newcommand{\tstep}{\mathrm{twostep}}
\newcommand{\unif}{\mathrm{uniform}}
\newcommand{\tv}{\mathrm{TV}}
\renewcommand{\ol}{\overline}

